
\section{Conocimientos, experiencia, tipo y autorización}
Parte 66 no se reduce a aprobar un examen. La ruta combina conocimientos, experiencia, categoría y, cuando corresponda, habilitación o anotación de tipo y condiciones para ejercer privilegios. 66.A.20 define privilegios y condiciones; .25 y .30, conocimientos y experiencia; .45, habilitaciones. Se lee el escenario concreto sin inventar una cifra de experiencia general para todos.

La autorización organizativa conecta persona y funciones dentro de un alcance. Una licencia no demuestra autorización para cualquier tarea y modelo; trabajar en un taller tampoco concede automáticamente privilegios de certificación. La evaluación de competencia y sus evidencias siguen los requisitos y procedimientos pertinentes.
\section{Formación y equivalencias demostradas}
Parte 147 regula organizaciones de formación en su ámbito. Reconocimiento de curso y certificados tiene condiciones. La universidad aporta bases científicas relevantes, pero el grado no se convierte por nombre en licencia. Posibles créditos o rutas requieren disposición concreta y evaluación competente.

Un expediente separa formación básica, formación de tipo, prácticas, experiencia y autorizaciones. Conserva fechas y contenidos sin inventar equivalencias. El laboratorio puede desarrollar habilidades valiosas con alcance diferente de la experiencia acreditada exigida por una ruta.
\section{Competencia continua y factores humanos}
Cambios de producto, datos y útiles requieren actualización. Familiaridad con una variante puede favorecer suposiciones incorrectas sobre otra. Trabajo reciente y condiciones de ejercicio se consideran en su régimen. Fatiga, comunicación, presión e interrupciones afectan ejecución.

Comprender el mecanismo material no sustituye el procedimiento; conocer una lista sin su propósito también puede generar errores ante una anomalía. La formación conecta conocimiento, ejecución y reconocimiento de límites.
\repaso{¿Carece de utilidad el grado para mantenimiento por no ser licencia?}{La base de análisis y materiales es relevante. Utilidad formativa y privilegio legal son conceptos distintos; se identifican los requisitos y evidencias del paso siguiente.}
